\documentclass[11pt,a4paper]{article}

% ------------------------------------------------------------------
%  Police et langue (compilation avec XeLaTeX)
% ------------------------------------------------------------------
\usepackage{fontspec}
\setmainfont{Latin Modern Roman}
\setsansfont{Latin Modern Sans}
\setmonofont{Latin Modern Mono}
\usepackage{polyglossia}
\setdefaultlanguage{french}
\usepackage{eurosym}

% ------------------------------------------------------------------
%  Mise en page
% ------------------------------------------------------------------
\usepackage[a4paper,margin=2.4cm,headheight=15pt]{geometry}
\usepackage[table]{xcolor}
\usepackage{titlesec}
\usepackage{enumitem}
\usepackage{tabularx}
\usepackage{array}
\usepackage{booktabs}
\usepackage{longtable}
\usepackage{fancyhdr}
\usepackage{lastpage}
\usepackage[hidelinks]{hyperref}

% ------------------------------------------------------------------
%  Couleurs de la charte
% ------------------------------------------------------------------
\definecolor{primary}{HTML}{1F4E79}
\definecolor{secondary}{HTML}{2E86AB}
\definecolor{light}{HTML}{EAF1F7}
\definecolor{greytext}{HTML}{555555}

\hypersetup{
  colorlinks=true,
  linkcolor=primary,
  urlcolor=secondary,
  pdftitle={Cahier des charges - Immobil'IA},
  pdfauthor={Equipe Immobil'IA}
}

% ------------------------------------------------------------------
%  Titres
% ------------------------------------------------------------------
\titleformat{\section}
  {\color{primary}\Large\bfseries}{\thesection.}{0.6em}{}
\titleformat{\subsection}
  {\color{secondary}\large\bfseries}{\thesubsection}{0.6em}{}
\titleformat{\subsubsection}
  {\color{primary}\normalsize\bfseries}{\thesubsubsection}{0.6em}{}
\titlespacing*{\section}{0pt}{18pt}{8pt}

% ------------------------------------------------------------------
%  En-têtes / pieds de page
% ------------------------------------------------------------------
\pagestyle{fancy}
\fancyhf{}
\renewcommand{\headrulewidth}{0.4pt}
\renewcommand{\footrulewidth}{0.4pt}
\fancyhead[L]{\small\color{greytext}Immobil'IA}
\fancyhead[R]{\small\color{greytext}Cahier des charges}
\fancyfoot[L]{\small\color{greytext}Document confidentiel}
\fancyfoot[R]{\small\color{greytext}Page \thepage\ / \pageref{LastPage}}

% ------------------------------------------------------------------
%  Environnement d'exigence
%  #1 = identifiant, #2 = intitulé
% ------------------------------------------------------------------
\newcommand{\exigence}[2]{%
  \par\vspace{6pt}%
  \noindent\colorbox{light}{%
    \parbox{\dimexpr\textwidth-2\fboxsep}{%
      \textbf{\color{primary}#1} \textemdash\ \textbf{#2}%
    }%
  }\par\vspace{2pt}%
}

\newenvironment{acceptation}
  {\vspace{2pt}\noindent\textit{Règles d'acceptation :}
   \begin{itemize}[leftmargin=1.4em,itemsep=1pt,topsep=2pt]}
  {\end{itemize}\vspace{4pt}}

\setlist[itemize]{label=\textcolor{secondary}{\textbullet}}

% ==================================================================
\begin{document}

% ================= PAGE DE GARDE =================
\begin{titlepage}
\centering
\vspace*{2.2cm}

{\color{primary}\rule{\textwidth}{2pt}}\\[0.5cm]
{\Huge\bfseries\color{primary} Cahier des charges}\\[0.35cm]
{\LARGE\color{secondary} Immobil'IA}\\[0.35cm]
{\large\color{greytext} Plateforme immobilière augmentée par l'intelligence artificielle}\\[0.3cm]
{\color{primary}\rule{\textwidth}{2pt}}\\[2cm]

\begin{tabular}{rl}
\textbf{Projet}  & : Immobil'IA \\[4pt]
\textbf{Type}    & : Cahier des charges fonctionnel \\[4pt]
\textbf{Version} & : 1.0 \\[4pt]
\textbf{Date}    & : \today \\[4pt]
\textbf{Statut}  & : À valider \\
\end{tabular}

\vfill
{\small\color{greytext} Document décrivant l'ensemble des fonctionnalités,
des critères et des règles d'acceptation attendus de la solution.}
\end{titlepage}

% ================= TABLE DES MATIÈRES =================
\tableofcontents
\thispagestyle{fancy}
\clearpage

% ==================================================================
\section{Présentation du projet}

\subsection{Contexte}
Immobil'IA est une plateforme immobilière qui met en relation les personnes
qui proposent un bien (vente ou location), celles qui recherchent un bien et
celles qui souhaitent investir. La plateforme s'appuie sur l'intelligence
artificielle pour assister chaque profil d'utilisateur dans ses tâches :
dépôt et optimisation d'annonces, recherche et comparaison de biens, analyse
d'investissement, ainsi que des services transverses d'aide à la décoration
et à l'estimation de travaux.

\subsection{Objectif du document}
Ce cahier des charges recense l'ensemble des fonctionnalités attendues de la
solution, regroupées par acteur, ainsi que les critères et les règles
d'acceptation associés à chacune d'elles. Il constitue le référentiel
fonctionnel partagé entre les parties prenantes du projet. Il ne préjuge
d'aucun choix technique, d'aucune architecture ni d'aucune technologie de
mise en œuvre.

\subsection{Périmètre}
Le périmètre couvre les fonctionnalités décrites pour les acteurs suivants :
le Vendeur / Bailleur, le Locataire / Acheteur, l'Investisseur, ainsi que les
fonctionnalités communes à tous les acteurs. Toute fonctionnalité non
mentionnée dans le présent document est considérée comme hors périmètre.

\subsection{Acteurs}
\begin{itemize}[leftmargin=1.4em]
  \item \textbf{Le Vendeur / Bailleur} : utilisateur qui dépose une offre de vente ou de location.
  \item \textbf{Le Locataire / Acheteur} : utilisateur qui recherche un bien à louer ou à acheter.
  \item \textbf{L'Investisseur} : utilisateur qui analyse des biens dans une optique de rentabilité et de placement.
  \item \textbf{Tous les acteurs} : fonctionnalités communes accessibles à l'ensemble des profils.
\end{itemize}

\subsection{Convention de lecture des exigences}
Chaque exigence fonctionnelle est identifiée par un code unique de la forme
\texttt{EF-X-NN}, où \texttt{X} désigne l'acteur concerné
(\texttt{V} : Vendeur / Bailleur, \texttt{L} : Locataire / Acheteur,
\texttt{I} : Investisseur, \texttt{T} : transverse / tous les acteurs) et
\texttt{NN} le numéro d'ordre. Chaque exigence est accompagnée d'une
description et d'une liste de règles d'acceptation. Une exigence n'est
considérée comme livrée que si l'ensemble de ses règles d'acceptation est
satisfait.

\clearpage

% ==================================================================
\section{Exigences fonctionnelles — Vendeur / Bailleur}

\exigence{EF-V-01}{Dépôt de l'offre}
Le vendeur / bailleur peut déposer une offre comprenant des photos, des
vidéos et une description complète du bien.
\begin{acceptation}
  \item L'utilisateur peut créer une nouvelle offre et y ajouter une ou plusieurs photos.
  \item L'utilisateur peut ajouter une ou plusieurs vidéos à l'offre.
  \item L'utilisateur peut saisir une description complète du bien.
  \item L'offre ne peut être publiée que si les éléments requis (photos, vidéos, description) sont fournis.
\end{acceptation}

\exigence{EF-V-02}{Recommandation de prix intelligente}
L'IA analyse l'offre et recommande un prix, en signalant tout écart par
rapport au marché.
\begin{acceptation}
  \item Le système propose un prix recommandé après analyse de l'offre déposée.
  \item Le système signale lorsque le prix saisi est trop élevé par rapport au marché.
  \item Le système signale lorsque le prix saisi est trop faible par rapport au marché.
  \item Le vendeur reste libre d'accepter ou d'ignorer la recommandation.
\end{acceptation}

\exigence{EF-V-03}{Recommandation des charges de location}
Le système recommande les charges de location à intégrer dans la description.
\begin{acceptation}
  \item Le système propose une recommandation de charges de location pour l'offre.
  \item La recommandation peut être intégrée à la description de l'annonce.
\end{acceptation}

\exigence{EF-V-04}{Analyse concurrentielle automatique}
Le système génère automatiquement une liste d'offres similaires afin de se
positionner face à la concurrence.
\begin{acceptation}
  \item Le système génère une liste d'offres similaires au bien déposé.
  \item Les offres similaires correspondent au même quartier et au même type de bien.
  \item La liste permet au vendeur de comparer son offre à la concurrence.
\end{acceptation}

\exigence{EF-V-05}{Vérification de la description}
Le système détecte les incohérences et les informations invalides ou
manquantes dans l'annonce.
\begin{acceptation}
  \item Le système détecte les incohérences entre la description et les photos (ex. surface, nombre de pièces).
  \item Le système signale les informations invalides.
  \item Le système signale les informations manquantes.
\end{acceptation}

\exigence{EF-V-06}{Génération de vidéo de visite virtuelle}
Le système génère une vidéo de visite virtuelle réaliste à partir des photos
et vidéos brutes fournies.
\begin{acceptation}
  \item Le système génère une vidéo de visite virtuelle à partir des photos et vidéos fournies par le vendeur.
  \item La vidéo générée est associée à l'offre.
\end{acceptation}

\exigence{EF-V-07}{Score de qualité de l'annonce}
Le système attribue un score de qualité à l'annonce, accompagné de
suggestions concrètes d'amélioration.
\begin{acceptation}
  \item Le système calcule et affiche un score de qualité de l'annonce.
  \item Le système fournit des suggestions concrètes d'amélioration (ex. angle de photo manquant, luminosité à corriger, description trop courte).
  \item Le score reflète la prise en compte des améliorations apportées.
\end{acceptation}

\exigence{EF-V-08}{Génération des réglementations et d'un contrat conforme}
Le système génère les réglementations applicables et un exemple de contrat
conforme à la loi française.
\begin{acceptation}
  \item Le système génère les réglementations applicables à l'offre.
  \item Le système génère un exemple de contrat conforme à la loi française.
\end{acceptation}

\exigence{EF-V-09}{Alerte anti-fraude / anti-discrimination}
Le système détecte les formulations discriminatoires ou illégales présentes
dans l'annonce.
\begin{acceptation}
  \item Le système détecte les formulations discriminatoires dans l'annonce.
  \item Le système détecte les formulations illégales dans l'annonce.
  \item Le système alerte le vendeur lorsqu'une telle formulation est détectée.
\end{acceptation}

\clearpage

% ==================================================================
\section{Exigences fonctionnelles — Locataire / Acheteur}

\exigence{EF-L-01}{Saisie du besoin en langage naturel}
Le locataire / acheteur exprime son besoin en langage naturel (budget,
critères, mode de vie), exploité via une recherche hybride.
\begin{acceptation}
  \item L'utilisateur peut saisir son besoin en langage naturel.
  \item La saisie prend en compte le budget, les critères et le mode de vie.
  \item La recherche s'appuie sur une approche de recherche hybride.
\end{acceptation}

\exigence{EF-L-02}{Agrégation d'offres}
Le système agrège les offres provenant de la base de la plateforme et du web.
\begin{acceptation}
  \item Le système présente les offres issues de la base de la plateforme.
  \item Le système présente les offres issues du web.
  \item Les offres des deux sources sont agrégées dans un même ensemble de résultats.
\end{acceptation}

\exigence{EF-L-03}{Cartographie comparative des offres}
Le système propose une visualisation géographique comparative des offres.
\begin{acceptation}
  \item Les offres sont affichées sur une représentation géographique.
  \item La visualisation présente le prix des offres.
  \item La visualisation présente les distances et les points d'intérêt (ex. bruit, pollution, embouteillage, école, université).
\end{acceptation}

\exigence{EF-L-04}{Recommandation des meilleures offres}
Le système recommande les meilleures offres selon le rapport qualité / prix
et une analyse concurrentielle par score.
\begin{acceptation}
  \item Le système recommande les offres selon le rapport qualité / prix.
  \item Les offres sont classées selon un score d'analyse concurrentielle (score d'honnêteté, par ordre croissant).
\end{acceptation}

\exigence{EF-L-05}{Alerte qualité de quartier}
Le système signale les quartiers mal notés.
\begin{acceptation}
  \item Le système signale lorsqu'un quartier est mal noté.
  \item Le signalement prend en compte des critères tels que la sécurité, le transport et le bruit.
\end{acceptation}

\exigence{EF-L-06}{Réglementations, contrat et vérification de contrat}
Le système génère les réglementations et un exemple de contrat conforme, et
vérifie le contrat généré par le vendeur / bailleur.
\begin{acceptation}
  \item Le système génère les réglementations applicables.
  \item Le système génère un exemple de contrat conforme à la loi française.
  \item Le système vérifie le contrat généré par le vendeur / bailleur.
\end{acceptation}

\exigence{EF-L-07}{Assistant de visite avec questions suggérées}
Le système prépare une liste de questions pertinentes à poser au propriétaire
selon le type de bien.
\begin{acceptation}
  \item Le système propose une liste de questions à poser au propriétaire.
  \item Les questions proposées sont adaptées au type de bien.
\end{acceptation}

\exigence{EF-L-08}{Chat comparatif multi-biens}
Le système compare plusieurs biens à la demande et génère un tableau de
synthèse.
\begin{acceptation}
  \item L'utilisateur peut demander la comparaison de plusieurs biens (ex. trois appartements).
  \item Le système génère un tableau de synthèse présentant les avantages et les inconvénients de chaque bien.
\end{acceptation}

\exigence{EF-L-09}{Matching de profils pour colocation}
Le système propose un appariement de profils pour la colocation lorsque
l'utilisateur en recherche.
\begin{acceptation}
  \item L'utilisateur peut indiquer qu'il recherche une colocation.
  \item Le système propose un appariement de profils correspondant à la recherche.
\end{acceptation}

\exigence{EF-L-10}{Vérification de disponibilité du bien}
Le système vérifie si le bien est déjà loué afin d'éviter les arnaques, via
une recherche par image.
\begin{acceptation}
  \item Le système permet de vérifier si un bien est déjà loué.
  \item La vérification peut s'appuyer sur une recherche par image.
\end{acceptation}

\exigence{EF-L-11}{Notation et commentaire}
L'utilisateur peut noter et commenter, ces éléments alimentant le score.
\begin{acceptation}
  \item L'utilisateur peut attribuer une note.
  \item L'utilisateur peut laisser un commentaire.
  \item La note contribue au calcul du score.
\end{acceptation}

\exigence{EF-L-12}{Signalement de scam / comportement inadéquat}
L'utilisateur peut signaler une arnaque ou un comportement inadéquat du
propriétaire.
\begin{acceptation}
  \item L'utilisateur peut signaler une arnaque (scam).
  \item L'utilisateur peut signaler un comportement inadéquat du propriétaire.
\end{acceptation}

\exigence{EF-L-13}{Sortie sous forme d'image du plan}
Le système peut produire une sortie sous forme d'image du plan.
\begin{acceptation}
  \item Le système peut générer une image représentant le plan du bien.
\end{acceptation}

\exigence{EF-L-14}{Notification d'offre similaire}
Si l'utilisateur n'a pas trouvé l'offre recherchée, il peut être notifié
lorsqu'une offre similaire est publiée.
\begin{acceptation}
  \item L'utilisateur peut demander à être notifié lorsqu'une offre similaire à sa recherche est publiée.
  \item Le système envoie une notification lors de la publication d'une offre similaire.
\end{acceptation}

\clearpage

% ==================================================================
\section{Exigences fonctionnelles — Investisseur}

\exigence{EF-I-01}{Rapport d'investissement personnalisé}
Le système génère un rapport d'investissement personnalisé avec simulation de
rentabilité locative.
\begin{acceptation}
  \item Le système génère un rapport d'investissement personnalisé.
  \item Le rapport inclut une simulation de rentabilité locative.
\end{acceptation}

\exigence{EF-I-02}{Scénarios « et si »}
Le système permet de simuler des scénarios « et si ».
\begin{acceptation}
  \item Le système permet de simuler l'impact de travaux.
  \item Le système permet de simuler l'évolution du quartier.
  \item Le système permet de simuler la variation des taux d'intérêt.
\end{acceptation}

\exigence{EF-I-03}{Score de risque locatif}
Le système calcule un score de risque locatif.
\begin{acceptation}
  \item Le système calcule et affiche un score de risque locatif.
  \item Le score prend en compte des facteurs tels que la zone tendue, la vacance locative moyenne et le profil des locataires types.
\end{acceptation}

\exigence{EF-I-04}{Comparateur multi-biens pour portefeuille}
Le système compare plusieurs biens pour orienter l'allocation d'un
portefeuille.
\begin{acceptation}
  \item L'utilisateur peut comparer plusieurs biens en vue d'un investissement.
  \item Le système aide à identifier le meilleur rendement pour un budget donné (ex. où investir 200\,k\euro{}).
\end{acceptation}

\exigence{EF-I-05}{Prévision de plus-value}
Le système fournit une prévision de plus-value à 5 et 10 ans basée sur les
tendances du quartier.
\begin{acceptation}
  \item Le système fournit une prévision de plus-value à 5 ans.
  \item Le système fournit une prévision de plus-value à 10 ans.
  \item La prévision se base sur les tendances du quartier.
\end{acceptation}

\exigence{EF-I-06}{Dashboard de l'investisseur}
Le système met à disposition un tableau de bord permettant à l'investisseur
de mieux comprendre ses intérêts.
\begin{acceptation}
  \item Le système affiche un tableau de bord dédié à l'investisseur.
  \item Le tableau de bord aide l'investisseur à comprendre ses intérêts.
\end{acceptation}

\clearpage

% ==================================================================
\section{Exigences fonctionnelles — Tous les acteurs}

\exigence{EF-T-01}{Génération d'idées de décoration}
À partir d'une photo ou d'une vidéo, le système génère des idées de
décoration selon l'activité et les préférences.
\begin{acceptation}
  \item L'utilisateur peut téléverser une photo ou une vidéo.
  \item Le système génère des idées de décoration à partir du média fourni.
  \item Les idées tiennent compte de l'activité et des préférences de l'utilisateur.
\end{acceptation}

\exigence{EF-T-02}{Décoration selon le budget et les préférences}
Le système propose une décoration selon le budget et les préférences, dans le
respect des réglementations.
\begin{acceptation}
  \item Les propositions de décoration respectent le budget indiqué.
  \item Les propositions de décoration respectent les préférences de l'utilisateur.
  \item Les propositions respectent les réglementations applicables.
\end{acceptation}

\exigence{EF-T-03}{Mode avant / après (home staging)}
Le système propose un mode comparatif « avant / après » pour convaincre un
vendeur de réaliser du home staging.
\begin{acceptation}
  \item Le système présente une comparaison « avant / après ».
  \item La comparaison illustre l'intérêt du home staging.
\end{acceptation}

\exigence{EF-T-04}{Estimation de réparation à partir des photos}
Le système estime les réparations à partir des photos et répercute cette
estimation sur le score et le prix.
\begin{acceptation}
  \item Le système détecte des besoins de réparation à partir des photos (ex. humidité, réparations dans les murs).
  \item L'estimation entraîne une diminution du score ou un ajustement de l'estimation de prix.
\end{acceptation}

\clearpage

% ==================================================================
\section{Récapitulatif des exigences}

\renewcommand{\arraystretch}{1.3}
\begin{longtable}{>{\bfseries}p{2.2cm} p{\dimexpr\textwidth-3cm}}
\toprule
\rowcolor{light}\textbf{Identifiant} & \textbf{Intitulé de l'exigence} \\
\midrule
\endfirsthead
\rowcolor{light}\textbf{Identifiant} & \textbf{Intitulé de l'exigence} \\
\midrule
\endhead
\bottomrule
\endfoot
EF-V-01 & Dépôt de l'offre \\
EF-V-02 & Recommandation de prix intelligente \\
EF-V-03 & Recommandation des charges de location \\
EF-V-04 & Analyse concurrentielle automatique \\
EF-V-05 & Vérification de la description \\
EF-V-06 & Génération de vidéo de visite virtuelle \\
EF-V-07 & Score de qualité de l'annonce \\
EF-V-08 & Génération des réglementations et d'un contrat conforme \\
EF-V-09 & Alerte anti-fraude / anti-discrimination \\
EF-L-01 & Saisie du besoin en langage naturel \\
EF-L-02 & Agrégation d'offres \\
EF-L-03 & Cartographie comparative des offres \\
EF-L-04 & Recommandation des meilleures offres \\
EF-L-05 & Alerte qualité de quartier \\
EF-L-06 & Réglementations, contrat et vérification de contrat \\
EF-L-07 & Assistant de visite avec questions suggérées \\
EF-L-08 & Chat comparatif multi-biens \\
EF-L-09 & Matching de profils pour colocation \\
EF-L-10 & Vérification de disponibilité du bien \\
EF-L-11 & Notation et commentaire \\
EF-L-12 & Signalement de scam / comportement inadéquat \\
EF-L-13 & Sortie sous forme d'image du plan \\
EF-L-14 & Notification d'offre similaire \\
EF-I-01 & Rapport d'investissement personnalisé \\
EF-I-02 & Scénarios « et si » \\
EF-I-03 & Score de risque locatif \\
EF-I-04 & Comparateur multi-biens pour portefeuille \\
EF-I-05 & Prévision de plus-value \\
EF-I-06 & Dashboard de l'investisseur \\
EF-T-01 & Génération d'idées de décoration \\
EF-T-02 & Décoration selon le budget et les préférences \\
EF-T-03 & Mode avant / après (home staging) \\
EF-T-04 & Estimation de réparation à partir des photos \\
\end{longtable}

\end{document}
